\scanpage{19}
\mspage{9}
Dualement, \deleted{\unsure{$A$ est de}} la donnée d'un carré-arêtes $A$ s'interprète comme la donnée d'un torseur sous un groupe $\mathbb{F}_2^D$, ce qui équivaut à la donnée d'une famille de \inserted{$\mathbb{F}_2$-}torseurs $(S_d)_{d\in D}$, \deleted{\unsure{indexée par un ens. $S$}}
\[
A\simeq\prod_{d\in D}S_d\qquad(\simeq\text{sections de $S$ sur $D$})
\]
-- la donnée de cette famille équivaut à la donnée d'un ensemble (à quatre éléments) $S$ et d'une application de degré $2$
\[
S\longrightarrow D
\]
qui avec la formule précédente donne
\[
A\hookrightarrow\mathcal{P}_2(S)
\]
et la relation d'incidence entre $S$ et $A$\ldots

\uline{Formulaire} On a, pour un carré $Q$, les trois ensembles à deux éléments
\[
\underset{\text{codiagonales}}{C},\quad
\underset{\text{diagonales}}{D},\quad
\underset{\text{orientations}}{\underline{\omega}}
\]
plus les deux paquets de deux ens. à deux éléments (ou $4$ autres ens. à deux éléments)
\[
\begin{array}{cc}
S\xrightarrow{\mathrm{deg.}\,2}D,&A\xrightarrow{\mathrm{deg.}\,2}C\\
(D\simeq S/\underline{a}_S).&(C\simeq A/\underline{a}_A)
\end{array}
\]
On a le formulaire d'iso. can.
% archive foliation: 18
